\phantomsection\section*{ВВЕДЕНИЕ}\addcontentsline{toc}{section}{ВВЕДЕНИЕ}

% База данных являются неотъемлемой частью любого программного обеспечения, где требуется хранить данные пользователя. 
% От базы данных напрямую зависит время выполнения запросов, испольующих сохраненные данные.
База данных является неотъемлемой частью любого программного обеспечения,
где требуется хранить данные пользователя.
От базы данных напрямую зависит время выполнения запросов, использующих сохраненные данные,
что критично для систем с высокой нагрузкой.
В системах управления проектами объем информации велик: необходимо ежедневно отслеживать статусы задач, сроки исполнения, приоритеты и взаимодействие между членами команды. Неоптимальная структура таблиц приводит к избыточности данных, нарушению ссылочной целостности и замедлению формирования аналитических отчетов. Поэтому грамотное проектирование схемы хранения становится ключевым фактором производительности всего приложения и обеспечения надежности бизнес-процессов.

Целью работы является разработка реляционной базы данных для системы управления проектами.

Для достижения поставленной цели необходимо решить следующие задачи: \\
\textbf{В аналитической части:}
\begin{enumerate}
\item Описать типовые процессы и роли в redmine, связанные с учебным процессом.
\item Описать сущности средства управления проектами redmine.
\item Проанализировать существующие решения в области СУБД.
\end{enumerate}
\textbf{В конструкторской части:}
\begin{enumerate}
\item Спроектировать базу данных.
\item Проверить, что разработанная база данных находится в 3НФ.
\item Разработать хранимые процедуру.
\item Разработать интеграции с внешним программным обеспечением для расширения функциональности.
\end{enumerate}
\textbf{В технологической части:}
\begin{enumerate}
    \item Обосновать выбор целевой платформы и средств разработки.
    \item Привести примеры реализации хранимых процедур, разработанных скриптов.
    \item Привести инструкции по развертыванию программы.
\end{enumerate}
\textbf{В исследовательской части:}
\begin{enumerate}
    \item Исследовать зависимость времени выполнения операций в зависимости от количества входящих запросов.
    \item Исследовать зависимость времени выполнения операций в зависимости от кэширования.
\end{enumerate}

\clearpage
